\emph{Setting.} The circulant reading of Proposition \claimref{synthesis.koide-channel-phase-is-no-symmetry-angle} reads the projected coherent vector $p_{k} = 1 + \sqrt2\cos(\delta + \tfrac{2\pi k}{3})$ as the $\sqrt{m}$ vector, that is, as the spectrum of a Hermitian circulant $C = a\mathbb{1} + bS + \bar{b}S^{\mathsf{T}}$, $b = |b|e^{i\delta}$; Proposition \claimref{synthesis.series-family-structures-give-no-lopsided-second-stage} places the channel phase on the coefficient of $A = S - S^{\mathsf{T}}$ in the same reading. Here the circulant is taken to act on the family basis in which the textures of Proposition \claimref{synthesis.theta23-at-s3-s2-direction} and Lemma \claimref{neutrino.breaking-relation} are written; without that identification the circulant is a list of numbers and no mixing statement follows. With $Sf_{k} = \omega^{k}f_{k}$ and $S^{\mathsf{T}}f_{k} = \omega^{-k}f_{k}$, $Cf_{k} = (a + 2|b|\cos(\delta + \tfrac{2\pi k}{3}))f_{k}$, so $p_{k}$ is the eigenvalue on $f_{k}$ and $f_{0} = u'/\sqrt3$ is the democratic line. At the stored representative ($0 < \delta < \tfrac{\pi}{12}$, $p_{1} < p_{2} < p_{0}$) the tau sits on $f_{0}$. A permutation $\sigma$ of the families acts by $C \mapsto P_{\sigma}CP_{\sigma}^{\mathsf{T}}$: a cyclic permutation fixes every circulant and a transposition sends $S$ to $S^{\mathsf{T}}$ ($\delta \mapsto -\delta$, exchanging $f_{1}$ and $f_{2}$). Neither moves $f_{0}$. The shift $\delta \mapsto \delta + \tfrac{2\pi}{3}$, which the circulant-reading result counts among relabellings of the $\sqrt{m}$ vector, sends $p_{k} \mapsto p_{k+1}$ and so changes which lepton sits on $f_{0}$; in a family basis shared with the neutrino sector it is a different matrix, so the placement on $f_{0}$ is a physical choice and all three placements are treated. Convention: here $(Sx)_{i} = x_{i+1}$, the transpose of the shift $Se_{j} = e_{j+1}$ of Proposition \claimref{synthesis.series-family-structures-give-no-lopsided-second-stage}, so $A = S - S^{\mathsf{T}}$ here is minus the $A$ there (the same line); clause (i) restates that Proposition's clauses (i) and (ii).

\emph{(i)} For any complex $c_{0}, c_{1}, c_{2}$, $(c_{0}\mathbb{1} + c_{1}S + c_{2}S^{2})f_{k} = (c_{0} + c_{1}\omega^{k} + c_{2}\omega^{2k})f_{k}$, and $CC^{\dagger}$ is again a polynomial in $S$; so both are diagonal on the $f_{k}$ whatever the channel phase. Under conjugation by $S_{3}$ the space of $3 \times 3$ matrices is $(1 \oplus \mathrm{std})^{\otimes 2} = 2 \cdot 1 \oplus \mathrm{sgn} \oplus 3 \cdot \mathrm{std}$. The circulants are the matrices fixed by conjugation with $S$; the $3$-cycle acts trivially on $1$ and $\mathrm{sgn}$ and with eigenvalues $\omega$, $\omega^{2}$ on $\mathrm{std}$, so the circulants are exactly the trivial and sign isotypic parts: the trivial part is $\mathrm{span}\{\mathbb{1}, \J\}$ and the sign part is the line of $A$, which changes sign under $P_{12}$ ($P_{12}SP_{12} = S^{\mathsf{T}}$). If $P_{12}MP_{12} = M$, then $\langle A, M\rangle = \langle P_{12}AP_{12}, P_{12}MP_{12}\rangle = -\langle A, M\rangle$, so $M$ has no component on $A$. A Hermitian circulant $\alpha\mathbb{1} + \beta S + \bar\beta S^{\mathsf{T}}$ ($\alpha$ real) is $\alpha\mathbb{1} + (\mathrm{Re}\,\beta)(\J - \mathbb{1}) + i(\mathrm{Im}\,\beta)A$, so it is $P_{12}$-invariant exactly when $\mathrm{Im}\,\beta = 0$, that is, when it lies in $\mathrm{span}\{\mathbb{1}, \J\}$, where $f_{1}$ and $f_{2}$ carry equal eigenvalues. The texture $(\J + \delta_{c}E)^{2} = 3\J + \delta_{c}(\J E + E\J) + \delta_{c}^{2}E$ is $P_{12}$-invariant; it is not a circulant for $\delta_{c} \ne 0$: a circulant has a constant diagonal, while the diagonal of $(\J + \delta_{c}E)^{2}$ is $(3, 3, 2 + (1 + \delta_{c})^{2})$, constant only at $\delta_{c} \in \{0, -2\}$, and at $\delta_{c} = -2$ the entries $(1,2)$, $(2,3)$, $(3,1)$ are $3, 1, 1$, not equal (the check confirms $S(\J + \delta_{c}E)^{2} \ne (\J + \delta_{c}E)^{2}S$ at $\delta_{c} = 0.286$). Both structures have the democratic $\J$ as leading term and the tau near $f_{0}$; beyond it they lie in different isotypic parts.

\emph{(ii)} If $m_{\nu} = e^{i\phi}m'$ with $m'$ real, the Takagi vectors (Fact \claimref{synthesis.imported-takagi-factorisation}) of $m_{\nu}$ are those of $m'$ times the phase $e^{-i\phi/2}$, which no modulus sees; so take $m_{\nu}$ real. For a real symmetric $m_{\nu}$ there is a real orthogonal $O$ with $m_{\nu} = O\,\mathrm{diag}(\lambda_{j})\,O^{\mathsf{T}}$, $\lambda_{j}$ real; the Takagi vectors are the columns of $O$ times phases ($1$ or $i$ to make the masses nonnegative), and with degenerate $\lambda_{j}$ a real basis can still be chosen. Moduli do not see the column phases. For a real vector $v$, $\langle f_{2}, v\rangle = \overline{\langle f_{1}, v\rangle}$, because $f_{2} = \bar{f_{1}}$; so $|U_{f_{1}j}| = |U_{f_{2}j}|$. And $\langle f_{0}, v\rangle = v\cdot u'/\sqrt3$. The check confirms the equality on ten exact real orthogonal frames (Cayley transforms of rational antisymmetric matrices, in $\mathbb{Q}(\sqrt{-3})$ arithmetic), that the overlaps form a unistochastic matrix, and, as a control, that a complex vector breaks it.

\emph{(iii)} For the textures of Proposition \claimref{synthesis.theta23-at-s3-s2-direction}, $\mD = d(\J + \delta E)$ and $M_{R} = \alpha\Pi + M_{3}\J/3$ commute with $P_{12}$ in the sense $P_{12}\mD P_{12} = \mD$, $P_{12}M_{R}P_{12} = M_{R}$, so $m_{\nu} = \mD M_{R}^{-1}\mD^{\mathsf{T}}$ is $P_{12}$-invariant, and $(1, -1, 0)$, the only odd vector up to scale, is an eigenvector; by clause (i) of that Proposition it is in the kernel. Adding any real $P_{12}$-invariant matrix to $E$ keeps the invariance and the eigenvector $(1, -1, 0)$, though not necessarily its zero eigenvalue. The check verifies the kernel and the invariance exactly at four rational instances (both roots of Proposition \claimref{synthesis.theta23-at-s3-s2-direction}(iv) with $\delta = \pm 0.286$, the peak and a generic point), and the eigenvector property with a rational $P_{12}$-invariant addition. This is the shared odd eigenvector of Proposition \claimref{synthesis.residual-s2-decouples-first-family-from-ckm-and-pmns}. Then $3|\langle f_{k}, (1, -1, 0)/\sqrt2\rangle|^{2} = |1 - \omega^{-k}|^{2}/2$ is $0$ for $k = 0$ and $\tfrac32$ for $k = 1, 2$, giving $(0, \tfrac12, \tfrac12)$; the even states $(a, a, b)$ give $(2a + b)^{2}/3$ on $f_{0}$ and $(a - b)^{2}/3$ on $f_{1}$ and $f_{2}$ (per unit norm). More generally, for any complex symmetric $m_{\nu}$ with $P_{12}m_{\nu}P_{12} = m_{\nu}$, write $v_{o} = (1, -1, 0)$, which spans the $P_{12}$-odd subspace. Then $P_{12}(m_{\nu}v_{o}) = m_{\nu}P_{12}v_{o} = -m_{\nu}v_{o}$, so $m_{\nu}v_{o} = \lambda v_{o}$ for some complex $\lambda$; since $v_{o}$ is real, $m_{\nu}\overline{(e^{i\psi}v_{o})} = \lambda e^{-i\psi}v_{o} = |\lambda|\,(e^{i\psi}v_{o})$ for $e^{2i\psi} = \lambda/|\lambda|$ (any $\psi$ if $\lambda = 0$), so $e^{i\psi}v_{o}/\sqrt2$ is a Takagi vector with singular value $|\lambda|$, and the other Takagi vectors can be chosen in the even subspace, which $m_{\nu}$ preserves. This covers the seesaw $m_{\nu} = \mD M_{R}^{-1}\mD^{\mathsf{T}}$ with $\mD = \alpha Y + \bar{\beta}\tilde{Y}$ as in Theorem \claimref{synthesis.rung-field-coupling-to-its-conjugate-is-the-second-family-matrix}, for complex frame coefficients and complex symmetric $P_{12}$-invariant $Y$, $\tilde{Y}$, $M_{R}$: the check builds six such instances in $\mathbb{Q}(\sqrt{-3})$ and verifies $P_{12}$-invariance and that $v_{o}$ is an eigenvector of $m_{\nu}$ and of $m_{\nu}m_{\nu}^{\dagger}$. That Theorem makes the Dirac neutrino matrix and $M_{R} = v_{R}f$ complex for generic couplings (real only on relatively real sub-branches), so this is the clause that reaches the covariant completion.

\emph{(iv)} In the standard parametrisation $|U_{e3}|^{2} = s_{13}^{2}$, $|U_{\mu 3}|^{2} = s_{23}^{2}c_{13}^{2}$, $|U_{\tau 3}|^{2} = c_{23}^{2}c_{13}^{2}$, $|U_{e1}|^{2} = c_{12}^{2}c_{13}^{2}$, $|U_{e2}|^{2} = s_{12}^{2}c_{13}^{2}$. With the ranges stated (from \cite{NuFIT2024}, Table 1, Fact \claimref{synthesis.imported-nufit-6-oscillation-parameter-ranges}; the union over normal and inverted ordering, with and without the Super-Kamiokande atmospheric data) the check computes in exact rationals $|U_{e3}|^{2} \in [0.02023, 0.02409]$, $|U_{e1}|^{2} \in [0.655 \times 0.97591, 0.725 \times 0.97977] \subset [0.639, 0.711]$, $|U_{e2}|^{2} \in [0.275 \times 0.97591, 0.345 \times 0.97977] \subset [0.268, 0.339]$, $|U_{\mu 3}|^{2} \ge 0.430 \times 0.97591 > 0.419$ and $|U_{\tau 3}|^{2} \ge 0.403 \times 0.97591 > 0.393$. Combining one-dimensional ranges in this way can only enlarge the allowed set, which is conservative for an exclusion. (a) $P_{12}$-invariant $m_{\nu}$, real or complex: by (iii) the column of the odd state has squared moduli $(0, \tfrac12, \tfrac12)$ on $(f_{0}, f_{1}, f_{2})$, so the electron's entry in that column is $0$ (electron on $f_{0}$) or $\tfrac12$ (electron on $f_{1}$ or $f_{2}$), whichever mass the odd state carries; neither value lies in any of the three electron-row ranges, so every placement is excluded. (b) Real $m_{\nu}$ (up to a phase): with the tau on $f_{0}$, by (ii) the electron and muon rows have equal moduli, so $|U_{e3}| = |U_{\mu 3}|$, impossible; with the muon on $f_{0}$, $|U_{e3}| = |U_{\tau 3}|$, impossible. Both hold for every real $m_{\nu}$, so for every real second stage at any order. (c) Electron on $f_{0}$ with a real $m_{\nu}$ that breaks $P_{12}$: by (ii) $|U_{\mu j}| = |U_{\tau j}|$. From column $3$, $c_{13}^{2}(s_{23}^{2} - c_{23}^{2}) = 0$, so $s_{23}^{2} = \tfrac12$; and the identity $|U_{\mu 1}|^{2} - |U_{\tau 1}|^{2} = (s_{12}^{2} - c_{12}^{2}s_{13}^{2})(c_{23}^{2} - s_{23}^{2}) + 4s_{12}c_{12}s_{13}s_{23}c_{23}\cos\delta_{CP}$ (checked at rational points of the circles) then gives $s_{12}c_{12}s_{13}\cos\delta_{CP} = 0$. This is the mu-tau reflection pattern of Harrison and Scott (\cite{HarrisonScott2002}, Fact \claimref{synthesis.imported-harrison-scott-mu-tau-reflection}; their treatment of the circulant basis is in an earlier paper, cited from memory). (d) A complex $P_{12}$-breaking $m_{\nu}$ escapes both (ii) and (iii): the check's control shows that complex vectors break the equality of the conjugate rows, and no statement is made there.

\emph{Scope and free numbers.} No parameter is fitted: the inputs are the stored placement of the circulant reading, the stored textures and the cited ranges; the look-elsewhere count is the three placements on $f_{0}$. The result says nothing if the circulant is not a matrix in the neutrino textures' family basis, and nothing for a complex neutrino sector that breaks $P_{12}$ (the generic case of the $\su(2)_{R}$-covariant completion, or a complex coefficient of a $120$). It is transport-absent.
